% -*- mode: latex; coding: utf-8 -*-
%-----------------------------------------------------------------------------
% pxGDX -- Disseny gràfic i portada per a P3CTeX
% Part of P3CTeX. Author: Scvria.
%
% Build from directory tex/ (TEXINPUTS so package is found):
%   set TEXINPUTS=.;./latex;./code;
%   pdflatex -interaction=nonstopmode doc/pxGDX.tex
%
% Or from tex/doc/: set TEXINPUTS=../latex:../code; then
%   pdflatex -interaction=nonstopmode pxGDX.tex
%-----------------------------------------------------------------------------
\documentclass[11pt,a4paper]{article}
\usepackage[T1]{fontenc}
\usepackage[utf8]{inputenc}
\usepackage{lmodern}
\usepackage{amsmath}
\usepackage{doc}
\usepackage[colorlinks,urlcolor=blue,linkcolor=blue,hyperindex]{hyperref}

\EnableCrossrefs

\title{El mòdul \texttt{pxGDX}\\
  \large Portada, estil de pàgina i metadades PDF a P3CTeX}
\author{Scvria\\P3CTeX / UOC PEC Repository}
\date{2026/02 -- Versió 0.1}

\begin{document}
\maketitle
\begin{abstract}
  El mòdul \texttt{pxGDX} encapsula la configuració gràfica de P3CTeX:
  portada, encapçalaments i peus de pàgina amb el format UOC, així com
  les metadades principals del fitxer PDF. Aquest document descriu la
  configuració recomanada a través de la classe \texttt{P3CTeX.cls} i
  les macros d'alt nivell disponibles.
\end{abstract}
\tableofcontents

\section{Introducció}

El codi intern de \texttt{pxGDX} viu al fitxer \verb|pxGDX.code.tex| i
es carrega automàticament per la classe \verb|P3CTeX.cls|. Els usuaris
no han d'invocar aquest fitxer directament: la manera recomanada és
sempre mitjançant la classe o, en casos especials, mitjançant les
macros d'alt nivell descrites en aquest manual.


\section{Activació i integració amb P3CTeX}

Quan es fa servir la classe \verb|P3CTeX.cls|, el mòdul \texttt{pxGDX}
es carrega de forma transparent i s'encarrega de:
\begin{itemize}
  \item definir la portada UOC;
  \item definir l'estil de pàgina (\emph{header/footer}) per a les
        pàgines normals i per al sumari;
  \item configurar els paràmetres bàsics de \texttt{hyperref} i les
        metadades del PDF.
\end{itemize}

La macro principal per configurar \texttt{pxGDX} a través de la classe
és:
\begin{verbatim}
\PECTeXconfig{<opcions>}
\end{verbatim}

On \verb|<opcions>| és un conjunt de claus que es distribueixen entre
els grups interns \verb|cover-config|, \verb|page-config| i
\verb|pdf-config|.

\section{Configuració de portada}

Les claus de \verb|cover-config| controlen el contingut i estil de la
portada. Un exemple típic:
\begin{verbatim}
\PECTeXconfig{
  cover-config = {
    text-backtitle = {Enginyeria Informàtica},
    text-uppertitle = {PAC 1},
    text-title = {Assignatura d'Exemple},
    text-minititle = {PAC1},
    text-author = {Nom de l'Estudiant}
  }
}
\end{verbatim}

Les claus més importants són:
\begin{itemize}
  \item \verb|text-backtitle|: títol de fons (grau o màster);
  \item \verb|text-title|: nom complet de l'assignatura;
  \item \verb|text-minititle|: codi o nom curt del document (per
        exemple, \verb|PAC1|);
  \item \verb|text-author|: nom de l'estudiant;
  \item \verb|text-date|: data que es mostrarà a la portada (per
        defecte, \verb|\today|).
\end{itemize}

Altres claus com \verb|uppertitle|, \verb|subtitle|, \verb|subsubtitle|
o \verb|add-picture| permeten activar o desactivar elements opcionals
de la portada. A continuació es detalla la llista completa de claus
de la família interna \verb|gdxCOVER|, amb els valors per defecte:

\subsection*{Claus de disseny (\texttt{gdxCOVER})}

\begin{description}
  \item{\verb|cover-header-height| ($\backslash$dim)} alçada del bloc superior de la
    portada (capçalera). Per defecte \verb|0.25\textheight|.
  \item{\verb|cover-subheader-height| ($\backslash$dim)} alçada del subbloc on es
    mostra el títol principal. Per defecte \verb|0.20\textheight|.
  \item{\verb|cover-body-height| ($\backslash$dim)} alçada del cos central de la
    portada (imatge i text auxiliar). Per defecte
    \verb|0.75\textheight|.
  \item{\verb|cover-square-size| ($\backslash$dim)} mida del requadre quadrat de la
    imatge central. Per defecte \verb|0.48\textwidth|.
  \item{\verb|cover-em| ($\backslash$dim)} unitat bàsica per a desplaçaments i
    marges interns. Per defecte \verb|0.08\textwidth|.
  \item{\verb|cover-font-max| ($\backslash$dim)} mida màxima de la tipografia de
    títols escalats. Per defecte \verb|42pt|.
\end{description}

\subsection*{Interruptors de contingut}

\begin{description}
  \item{\verb|uppertitle| (bool)} mostra el títol superior opcional
    (per exemple, ``PAC 1'') damunt del codi curt. Per defecte es
    considera \emph{fals} però la meta-clau \verb|fullcover| l'activa.
  \item{\verb|subtitle| (bool)} activa el subtítol sota el títol
    principal.
  \item{\verb|subsubtitle| (bool)} activa una franja amb informació
    addicional (normalment el codi de l'assignatura o un text curt).
  \item{\verb|add-picture| (bool)} activa la inserció de la imatge
    principal de portada (\verb|set-mainpicture| / \verb|set-auxpicture|).
  \item{\verb|outter-backtitle| (bool)} dibuixa el títol de fons sobre
    una franja de color a la part superior de la pàgina.
  \item{\verb|inner-backtitle| (bool)} repeteix el títol de fons dins
    del requadre principal.
  \item{\verb|noplagi| (bool)} controla si s'inclou el text de
    declaració de no plagi al peu de la portada.
  \item{\verb|offset| (bool)} aplica un petit desplaçament (\emph{offset})
    per aconseguir un efecte de superposició entre blocs.
  \item{\verb|clip-author| (bool)} controla si el nom de l'autor es
    mostra amb un efecte de ``retall'' de color.
  \item{\verb|unlink| (bool)} desactiva certs enllaços interns de
    \texttt{hyperref} a la portada (per evitar sobrecarregar la
    pàgina).
\end{description}

\subsection*{Textos i imatges de portada}

\begin{description}
  \item{\verb|text-backtitle|} títol de fons (grau, màster, etc.).
  \item{\verb|text-uppertitle|} títol superior opcional.
  \item{\verb|text-title|} títol principal (nom complet de
    l'assignatura).
  \item{\verb|text-minititle|} nom curt del document (PAC, PR, etc.).
  \item{\verb|text-subtitle|} subtítol opcional del document.
  \item{\verb|text-subsubtitle|} línia addicional (codi de
    l'assignatura, per exemple).
  \item{\verb|text-author|} nom de l'estudiant.
  \item{\verb|text-date|} data de lliurament; per defecte \verb|\today|.
  \item{\verb|text-auxtext|} text auxiliar a la part inferior del bloc
    principal; per defecte indica que el document s'ha fet amb P3CTeX.
  \item{\verb|set-mainpicture|} nom base de la imatge vertical
    principal (sense \verb|img/|); per defecte \verb|pxlngcover|.
  \item{\verb|set-auxpicture|} nom base de la imatge quadrada secundària;
    per defecte \verb|pxsqrcover|.
\end{description}

\subsection*{Colors, tipografia i marges}

\begin{description}
  \item{\verb|maincolor|} color principal de la portada (franges i
    elements destacats); per defecte \verb|secondaryColor|.
  \item{\verb|auxcolor|} color auxiliar per a caixes i fons; per
    defecte \verb|lightColor|.
  \item{\verb|textcolor|} color principal del text.
  \item{\verb|textcolor-aux|} color auxiliar del text sobre fons
    foscos; per defecte \verb|lighterColor|.
  \item{\verb|bg-color|} color de fons de les imatges i blocs; per
    defecte \verb|lighterColor|.
  \item{\verb|title-size|} macro de mida per al títol principal
    (per defecte \verb|\Huge|).
  \item{\verb|subtitle-size|} macro de mida per al subtítol
    (per defecte \verb|\large|).
  \item{\verb|margin-h| ($\backslash$dim)} marge horitzontal de la portada;
    per defecte \verb|24pt|.
  \item{\verb|margin-v| ($\backslash$dim)} marge vertical de la portada;
    per defecte \verb|36pt|.
\end{description}

\subsection*{Meta-claus}

\begin{description}
  \item{\verb|fullcover|} activa un conjunt coherent de booleans per
    obtenir una portada ``completa'' (tots els elements visibles).
  \item{\verb|print-placeholders|} configura la portada perquè mostri
    textos de mostra (placeholders) per a totes les etiquetes, útil
    per a depuració o com a guia per emplenar metadades.
\end{description}

\section{Configuració de la pàgina}

Les claus de \verb|page-config| controlen el text dels encapçalaments
i peus de pàgina (\emph{layout} UOC). Per exemple:
\begin{verbatim}
\PECTeXconfig{
  page-config = {
    author-name = {Nom de l'Estudiant},
    doc-name    = {PAC 1},
    subj-name   = {Nom Curt Assignatura},
    refcode     = {01312},
    course-name = {Enginyeria Informàtica}
  }
}
\end{verbatim}

Aquests valors es mostren a les capçaleres i als peus de pàgina de
les pàgines habituals i del sumari, mantenint la identitat visual
coherent en tot el document.

\subsection*{Claus de \texttt{gdxSTYLE}}

La família \verb|gdxSTYLE| agrupa totes les claus relacionades amb el
text i la disposició dels encapçalaments i peus:

\begin{description}
  \item{\verb|author-name|} nom de l'estudiant que es mostra al
    capçal de la pàgina (i que també s'utilitza com a valor per
    defecte de \verb|pdf-author|).
  \item{\verb|doc-name|} nom curt del document (PAC, PR, etc.) que
    apareix al centre de la capçalera.
  \item{\verb|subj-name|} nom curt de l'assignatura.
  \item{\verb|refcode|} codi de l'assignatura o referència acadèmica.
  \item{\verb|course-name|} nom complet del grau o màster.
  \item{\verb|header-left|} codi \LaTeX\ per al bloc esquerre de la
    capçalera; per defecte inclou el logotip UOC.
  \item{\verb|header-center|} codi del bloc central de la capçalera;
    per defecte mostra el nom del document i l'estudiant.
  \item{\verb|header-right|} codi del bloc dret; per defecte mostra el
    nom curt de l'assignatura i el codi.
  \item{\verb|footer-left|} codi del peu esquerre (buit per defecte).
  \item{\verb|footer-center|} codi del peu central; per defecte mostra
    la numeració ``pàgina X de Y''.
  \item{\verb|footer-right|} codi del peu dret; per defecte inclou el
    logotip secundari.
\end{description}

\section{Metadades PDF}

Finalment, les claus de \verb|pdf-config| permeten establir les
metadades principals del fitxer PDF generat a través de
\texttt{hyperref}:
\begin{verbatim}
\PECTeXconfig{
  pdf-config = {
    pdf-title  = {Títol del Document},
    pdf-author = {Nom de l'Estudiant},
    pdfsubject = {Assignatura o àrea},
    pdfkeywords = {P3CTeX, UOC, LaTeX}
  }
}
\end{verbatim}

Les claus \verb|linkcolor|, \verb|citecolor| i \verb|urlcolor| es poden
ajustar si cal canviar els colors per defecte dels enllaços.

\subsection*{Claus de \texttt{gdxPDF} i \texttt{pxGDX}}

Internament, \verb|pdf-config| es correspon amb la família de claus
\verb|gdxPDF|, que controla totes les metadades passades a
\texttt{hyperref}:

\begin{description}
  \item{\verb|pdf-title|} títol del document (camp \texttt{/Title} del
    PDF).
  \item{\verb|pdf-author|} autor del document; per defecte pren el
    valor d'\verb|author-name|.
  \item{\verb|pdfsubject|} assumpte o descripció breu del contingut.
  \item{\verb|pdfkeywords|} paraules clau separades per comes.
  \item{\verb|pdf-creator|} cadena lliure que identifica l'aplicació
    que ha generat el fitxer.
  \item{\verb|pdf-producer|} producte utilitzat com a ``producer'' del
    PDF; per defecte \verb|P3C.TeX v.3|.
  \item{\verb|linkcolor|} color dels enllaços interns.
  \item{\verb|citecolor|} color de les cites.
  \item{\verb|urlcolor|} color dels enllaços URL.
  \item{\verb|filecolor|} color dels enllaços a fitxers.
  \item{\verb|bookmarks| (bool)} activa o desactiva els marcadors de
    navegació (bookmarks) al PDF.
\end{description}

A nivell més alt, la família \verb|pxGDX| ofereix les claus següents,
que són les que utilitza \verb|\PECTeXconfig|:

\begin{description}
  \item{\verb|cover-config|} aplega totes les claus de \verb|gdxCOVER|
    per configurar la portada.
  \item{\verb|page-config|} aplega les claus de \verb|gdxSTYLE| per a
    encapçalaments i peus.
  \item{\verb|pdf-config|} aplega les claus de \verb|gdxPDF| per a les
    metadades del PDF.
  \item{\verb|src-path|} prefix de camí per a recursos gràfics
    utilitzats per \texttt{pxGDX} (per defecte buit).
\end{description}

\section{Macros d'ajuda per al cos del document}

A més de la configuració, \texttt{pxGDX} proporciona algunes macros
útils per al cos del document:

\begin{description}
  \item{\verb|\sectionuoc{<titol>}|} Crea una secció amb el color
    corporatiu UOC.
  \item{\verb|\subsectionuoc{<titol>}|} Variant per a subseccions.
  \item{\verb|\subsubsectionuoc{<titol>}|} Variant per a subsubseccions.
  \item{\verb|\pxHRule|} Dibuixa una línia horitzontal; la versió amb
    estrella i/o argument opcional p ermet controlar gruix i amplada.
  \item{\verb|\pxPageRule|} Dibuixa una línia horitzontal centrada a
    la pàgina, pensada per a separadors dins d'enunciats.
\end{description}

\section{Exemple mínim amb P3CTeX}

\begin{verbatim}
\documentclass[fake]{P3CTeX}

\begin{document}

\PECTeXconfig{
  cover-config = {
    text-minititle = {PAC1},
    text-author    = {Nom de l'Estudiant}
  },
  page-config = {
    doc-name   = {PAC 1},
    subj-name  = {Nom Curt},
    refcode    = {01312}
  }
}

\PECTeXprintCover*
\PECTeXprintTOC*

\sectionuoc{Introducció}
Text d'exemple amb l'estil UOC.

\pxPageRule*

\end{document}
\end{verbatim}

\section{Historial de canvis}

\begin{description}
  \item[v0.1] Versió inicial de la documentació de \texttt{pxGDX}:
    resum de la configuració de portada, pàgina i metadades, i llistat
    de macros d'ajuda principals.
\end{description}

\PrintIndex
\end{document}

